\section{Covariant structure: from \texorpdfstring{$\gamma^{\mu}$}{gamma\^mu} to \texorpdfstring{$SL(2,\mathbb{C})$}{SL(2,C)}}
\label{p:covariant}

At this point the discrete construction has reached a four-component
Dirac block whose squared dispersion is Lorentzian. The remaining task
is to identify the covariant structure that organizes it: a Clifford
algebra, the spinorial generators of Lorentz, and the connection to
$SL(2,\bbC)$.

\subsection{Clifford algebra}
In the Dirac representation, $\gamma^{0} = \tau_{z}\otimes\bbI$ and
$\gamma^{i} = \ii\,\tau_{x}\otimes\sigma_{i}$. A direct check confirms
\begin{equation}
\{\gamma^{\mu},\gamma^{\nu}\}\;=\;2\,\eta^{\mu\nu}\,\bbI_{4} ,
\qquad \eta=\mathrm{diag}(+1,-1,-1,-1) ,
\label{p:eq:clifford}
\end{equation}
which we verify entry by entry over all sixteen pairs $(\mu,\nu)$
(\texttt{validators/clifford.py}).

\subsection{Spinorial Lorentz generators}
The six spinorial Lorentz generators are
\begin{equation}
\Sigma^{\mu\nu} \;=\; \frac{\ii}{4}\,[\gamma^{\mu},\gamma^{\nu}] .
\end{equation}
Two structural facts then follow:
\begin{enumerate}[leftmargin=*]
\item the covariance identity
\begin{equation}
[\Sigma^{\rho\sigma},\,\gamma^{\mu}]
\;=\;\ii\bigl(\eta^{\sigma\mu}\gamma^{\rho}-\eta^{\rho\mu}\gamma^{\sigma}\bigr) ;
\label{p:eq:cov-id}
\end{equation}
\item the closure of the Lorentz algebra,
\begin{equation}
[\Sigma^{\mu\nu},\Sigma^{\rho\sigma}] = \ii\bigl(
  \eta^{\nu\rho}\Sigma^{\mu\sigma}-\eta^{\mu\rho}\Sigma^{\nu\sigma}
 -\eta^{\nu\sigma}\Sigma^{\mu\rho}+\eta^{\mu\sigma}\Sigma^{\nu\rho}\bigr).
\label{p:eq:lorentz-closure}
\end{equation}
\end{enumerate}
Both are verified over all 256 quadruples $(\mu,\nu,\rho,\sigma)$ by
\texttt{validators/lorentz.py::test\_sigma\_mu\_nu\_covariance} and
\texttt{::test\_lorentz\_algebra\_closure}. Rotational generators
$\Sigma^{ij}$ are Hermitian; boost generators $\Sigma^{0i}$ are
anti-Hermitian, consistent with rotations being compact and boosts being
non-compact.

\subsection{Double cover $SL(2,\bbC)\!\to\!SO^{+}(3,1)$}
A 2$\pi$ rotation in physical space acts on a spinor as
$\mathrm{e}^{-\ii\pi\sigma_{z}} = -\bbI$, not $+\bbI$, while a 4$\pi$
rotation returns the spinor to itself. This is the standard double-cover
relation, and is verified directly by
\texttt{validators/puentes\_grupos.py::test\_two\_pi\_gives\_minus\_identity}
and \texttt{::test\_four\_pi\_gives\_identity}. The chiral matrix
$\gamma_{5} = \ii\,\gamma^{0}\gamma^{1}\gamma^{2}\gamma^{3}$ anticommutes
with each $\gamma^{\mu}$ and squares to $\bbI_{4}$, defining the chiral
projectors $P_{\pm}=(\bbI\pm\gamma_{5})/2$ that decompose the four-spinor
into Weyl pairs --- closing the loop with the splitting of
Sec.~\ref{p:weyl-dirac-3d}.
